\section{Guía de instalación}

Dependencias:
\begin{itemize}[nosep]
  \item compilador C11 (\texttt{gcc} o \texttt{clang});
  \item \texttt{make};
  \item \texttt{python3} para integraciones y stress;
  \item \texttt{valgrind} en Linux/Pampero para memoria/fds;
  \item \texttt{matplotlib} si se quieren regenerar gráficos de stress;
  \item \texttt{rsync} y \texttt{ssh} para el runner de Pampero.
\end{itemize}

Compilación:
\begin{verbatim}
make clean && make
\end{verbatim}

Artefactos generados:
\begin{itemize}[nosep]
  \item \texttt{bin/server}: servidor SOCKS5 + PMC;
  \item \texttt{bin/client}: cliente CLI del PMC;
  \item \texttt{docs/report/main.pdf}: este informe.
\end{itemize}

Verificación remota recomendada:
\begin{verbatim}
PAMPERO_USER=<usuario> bash scripts/run-on-pampero.sh 12080
\end{verbatim}

El script sube el repo a Pampero, recompila desde cero y ejecuta unitarios,
integraciones sobre sockets y Valgrind con tráfico real.
